\documentclass[final]{beamer}
\usepackage[orientation=portrait,size=custom,width=90,height=130,scale=1.15]{beamerposter}
\usepackage{lmodern}
\renewcommand{\familydefault}{\sfdefault}
\usepackage[T1]{fontenc}
\usepackage{graphicx,booktabs,amsmath,amssymb,array}
\usepackage{dancheong}
\graphicspath{{figures/}}
\title{Persistent Homology of Music Network with Three Different Distances}
\begin{document}
\begin{frame}[t]
% ===================== HEADER =====================
\begin{beamercolorbox}[wd=\paperwidth,sep=1.2ex,leftskip=2ex,rightskip=2ex]{headline}
\begin{columns}[c]
  \begin{column}{0.13\paperwidth}\centering
    % institutional logos (drop real logos into poster/figures/ later; text fallback now)
    {\color{dcPaper}\large POSTECH\\MINDS\\KIAS}
  \end{column}
  \begin{column}{0.62\paperwidth}\centering
    {\color{dcPaper}\veryHuge\bfseries Persistent Homology of Music Network}\\[0.4ex]
    {\color{dcPaper}\Huge\bfseries with Three Different Distances}\\[0.8ex]
    {\color{dcPaper}\large \underline{Eunwoo Heo}$^{1,2,\ast}$ \quad Byeongchan Choi$^{1,2}$ \quad
     Myungock Kim$^{3}$ \quad Mai Lan Tran$^{1,2}$ \quad Jae-Hun Jung$^{1,2}$}\\[0.4ex]
    {\color{dcPaper}\normalsize $^{1}$Dept.\ of Mathematics, POSTECH \quad
     $^{2}$POSTECH MINDS \quad $^{3}$KIAS \quad
     {\itshape $^{\ast}$Present address: UNIST, Ulsan, Korea}}
  \end{column}
  \begin{column}{0.13\paperwidth}\centering
    
\includegraphics[width=0.7\linewidth]{qr_project.png}\\[0.2ex]
    {\color{dcPaper}\small Project page}
  \end{column}
\end{columns}
\end{beamercolorbox}
\vskip1.2ex

% ===================== CONTENT =====================
% ---------- HERO STRIP ----------
\begin{columns}[t]
  \begin{column}{0.40\paperwidth}
    \begin{block}{The big idea}
      \large
      Distance is a \textbf{tunable resolution knob} on musical structure.
      Sweeping \done${}\to{}$\dthree${}\to{}$\dtwo\ distills a piece down to its
      \textbf{core cyclic skeleton}.\par\vskip0.6ex
      \normalsize
      \begin{itemize}
        \item Three path distances on the music graph: \done\ (min edges),
              \dtwo\ (min weight, the only metric), \dthree\ (hybrid).
        \item Proven ordering \;$\textcolor{dcD2}{d_2}\!\le\!\textcolor{dcD3}{d_3}\!\le\!\textcolor{dcD1}{d_1}$\;
              and barcode injection $B_2\!\subseteq\!B_3\!\subseteq\!B_1$.
        \item Validated on real Korean court music (exact match to the paper's results).
      \end{itemize}
    \end{block}
  \end{column}
  \begin{column}{0.59\paperwidth}
    \begin{block}{Cycle skeleton extraction \;\textmd{(Sanghyeondodeuri, geomungo)}}
      \centering
      \includegraphics[width=\linewidth]{signature.pdf}\\[0.3ex]
      {\normalsize As \done${}\to{}$\dthree${}\to{}$\dtwo, cycles vanish
      \textbf{7$\,\to\,$4$\,\to\,$3}: only the robust core loops survive.}
    \end{block}
  \end{column}
\end{columns}
\vskip0.6ex
% ---------- 3 CONTENT COLUMNS ----------
\begin{columns}[t]
  \begin{column}{0.32\paperwidth}
    % === COLUMN 1 (Task 4) ===
  \end{column}
  \begin{column}{0.32\paperwidth}
    % === COLUMN 2 (Task 5) ===
  \end{column}
  \begin{column}{0.32\paperwidth}
    % === COLUMN 3 (Task 6) ===
  \end{column}
\end{columns}

\vfill
% ===================== FOOTER =====================
\begin{beamercolorbox}[wd=\paperwidth,sep=0.8ex,leftskip=2ex,rightskip=2ex]{headline}
\begin{columns}[c]
  \begin{column}{0.7\paperwidth}
    {\color{dcPaper}\small Key refs: Heo--Choi--Jung (2024) arXiv:2501.03553; Tran--Park--Jung (2023) \emph{J.\ Math.\ Music}; Sethares--Budney (2014).
    \quad Funded by NRF Korea (2021R1A2C3009648) \& KIAS Transdisciplinary program.
    \quad Contact: hew0920@postech.ac.kr}
  \end{column}
  \begin{column}{0.28\paperwidth}\raggedleft
    
\includegraphics[height=4cm]{qr_arxiv.png}\;{\color{dcPaper}\small arXiv:2506.13595}
  \end{column}
\end{columns}
\end{beamercolorbox}
\end{frame}
\end{document}
